\section{Experiments}

\subsection{Experimental Setup}

All experiments were conducted on an Apple Silicon Mac (MPS backend) using PyTorch 2.8.0, torchvision 0.23.0, and NumPy 2.0.2. The random seed was fixed to 42 for the current reproducibility run; random number generators in Python, NumPy, and PyTorch were seeded, and deterministic MPS operations were enabled where available. Results are from a single seed and should not be interpreted as a multi-seed estimate. Due to floating-point non-determinism in the MPS backend, exact numerical reproduction on different hardware is not guaranteed.

Hyperparameters for Stages 2 and 3 are held identical (lr\,$=$\,$5{\times}10^{-5}$, batch\,$=$\,32, weight decay\,$=$\,$10^{-4}$, pos\_weight\,$=$\,dynamic from training split) to enable valid causal comparison. Stage 1 uses different hyperparameters as a pretrained initialization.

\subsection{Baseline}

The baseline model is trained exclusively on HAM10000 (7,818 mel/nv samples) with an 80/20 stratified split. No DDI data is used at this stage. The model achieved validation AUROC 0.9650 on HAM10000, and training used the documented 10-epoch limit with early stopping.

When evaluated on the DDI test set using a threshold selected by Youden's $J$ on the validation set, the baseline achieved Dark AUROC\,$=$\,0.588 and sensitivity\,$=$\,0.400. Light-skin AUROC was 0.719 and overall test AUROC was 0.662. These results show a skin-tone performance gap, consistent with documented disparities in dermatological AI systems \cite{Daneshjou2022, Groh2021}.

\subsection{Fine-tuning}

Fine-tuning the baseline model on DDI (393 training samples) reduced overall test AUROC from 0.662 to 0.634. Light AUROC was 0.756 and Dark AUROC was 0.547, changing the subgroup AUROC gap from $-0.131$ to $-0.209$ (Dark minus Light). At the validation-selected threshold, Light and Dark sensitivities were both 0.600.

This result demonstrates that standard fine-tuning on a small, skin-tone-imbalanced dataset can \emph{amplify rather than mitigate} the performance gap. The fine-tuned model learns to exploit the overrepresented Light subgroup features while failing to generalize to Dark subgroup patterns.

Bootstrap 95\% confidence intervals for the fine-tuned model are: Light AUROC\,$=$\,0.761\,[0.575--0.903], Dark AUROC\,$=$\,0.548\,[0.341--0.750]. The wide Dark CI reflects the small subgroup size ($n\,{=}\,42$, 10 melanomas).

\subsection{Synthetic Augmentation}

Adding 290 synthetic dark-skin melanoma images reduced overall test AUROC slightly from 0.634 to 0.640 relative to the fine-tuned checkpoint, but reduced Dark AUROC from 0.547 to 0.481. The Light--Dark AUROC gap widened to $-0.291$. The paired Dark AUROC delta was $-0.066$ with 95\% bootstrap interval $[-0.151,0.008]$ and $p=0.962$. Sensitivity was 0.600 in both Light and Dark subgroups at the stage-specific validation-selected threshold.

Bootstrap 95\% CIs are: Light AUROC\,$=$\,0.775\,[0.582--0.925], Dark AUROC\,$=$\,0.481\,[0.276--0.688]. With only 10 melanomas per subgroup, the study is underpowered. Multi-seed repetition, external validation, and clinically meaningful synthetic-image review remain necessary.

As noted in Section~\ref{sec:synthetic_aug}, synthetic images are now generated exclusively from training-split data, eliminating the information leakage present in earlier versions of this work.

\subsection{Ablation Study}

We conduct an ablation study over the number of saved synthetic images, testing multipliers of 0$\times$ (no synthetics), 2$\times$ (58 images), 5$\times$ (145 images), and 10$\times$ (290 images). All synthetic images are generated from training-split dark-skin melanomas. Starting from the fine-tuned checkpoint, each variant is trained for up to 5 epochs with identical hyperparameters to Stage 3.

\paragraph{Methodology.}
Unlike earlier work that tested on-the-fly oversampling (duplicating real images with random transforms during training), this ablation uses \emph{saved synthetic images} generated via fixed Albumentations transforms. This approach is methodologically consistent with the main pipeline and enables direct comparison of augmentation quantity effects.

Results are summarized in Table~\ref{tab:ablation} and visualized in Figure~\ref{fig:ablation_tradeoff}. The ablation reveals a non-monotonic relationship between synthetic data quantity and Dark subgroup performance: Dark AUROC was 0.5469 at 0$\times$, 0.5750 at 2$\times$, 0.5344 at 5$\times$, and 0.4812 at 10$\times$. Overall test AUROC was 0.6848, 0.7063, 0.6742, and 0.6433, respectively. Validation AUROC peaked at 2$\times$, so no single multiplier was uniformly optimal.

Dark-subgroup AUROC peaked at 2$\times$ (0.575) and was lowest at 10$\times$ (0.481). This is an observed single-seed pattern, not proof of overfitting or synthetic artifacts; multi-seed repetition is needed to establish whether it is stable.
